\begin{table}[ht]
  \centering
  \caption{Three-Channel Decomposition: Urgent vs.\ Ordinary Procurement}
  \label{tab:three_channel_uo}
  \small
  \begin{threeparttable}
  \begin{tabular}{lcccc}
    \hline\hline
                                    & (1)         & (2)            & (3)        & (4)             \\
                                    & Baseline    & + log qty.\   & + firm FE  & + qty + firm FE \\
    \hline
    Urgent purchase & 0.051*** & 0.028 & 0.024 & 0.011 \\
                                    & (0.015) & (0.018) & (0.016) & (0.018) \\[3pt]
    Implied premium (\%)            & 5.26\% & 2.85\% & 2.48\% & 1.10\% \\
    \hline
    Item FE                          & Yes & Yes & Yes & Yes \\
    Year FE                          & Yes & Yes & Yes & Yes \\
    PBU FE                           & Yes & Yes & Yes & Yes \\
    Firm FE                          & No  & No  & Yes & Yes \\
    log Quantity control             & No  & Yes & No  & Yes \\
    \hline
    Observations                     & 196,883 & 196,883 & 196,642 & 196,642 \\
    \hline\hline
  \end{tabular}
  \begin{tablenotes}
    \small
    \item \textit{Notes:} Three-channel decomposition. Column (1) gives
    the total premium under the preferred specification. Columns (2)--(4)
    progressively absorb channels: (2) absorbs \emph{C1 demand fragmentation}
    via log quantity; (3) absorbs \emph{C3 supplier composition shift}
    via firm fixed effects; (4) absorbs both, leaving the C2
    \emph{competition channel} residual (and any within-firm markup).
    Decomposition: total 5.26\% = 2.41 pp (C1 quantity) + 2.79 pp (C3 firm
    selection) + 1.10 pp (residual).
    Standard errors clustered at the PBU level in parentheses.
    *** $p<0.01$, ** $p<0.05$, * $p<0.10$.
  \end{tablenotes}
  \end{threeparttable}
\end{table}
